\section{Procedury reakcji i przywrócenia}
\subsection*{Pożar lub uruchomienie ROP}
Wyzwalacz: dowolny alarm CSP. Odpowiedzialny: dyżurny SOC; patrol prowadzi ewakuację i udostępnia dojazd. SSP alarmuje bez oczekiwania na ocenę obrazu. Dyżurny odczytuje adres z CSP, niezwłocznie wzywa pomoc pod 112, kieruje personel do wyjść i pomocy osobom wymagającym wsparcia. Nie uzależnia reakcji od widocznego dymu w CCTV. PWP odłącza zasilanie niepożarowe wraz z wyjściem UPS zgodnie z jego fizycznym interfejsem. Przywrócenie: po potwierdzeniu bezpieczeństwa przez kierującego działaniami, usunięciu przyczyny i sprawdzeniu sterowań; reset ROP kluczem i CSP lokalnie na poziomie 2. Nie kasuje się alarmu z komputera VMS.

\subsection*{Włamanie, napad lub sabotaż}
Wyzwalacz: Z1--Z10 w odpowiednim czuwaniu, Z13/Z14 sabotaż lub Z15 napad całodobowo. Odpowiedzialny: dyżurny i patrol. Dyżurny odczytuje linię z KP-I01, wybiera odpowiednią kamerę i kieruje patrol bez narażania go na bezpośredni kontakt z napastnikiem; w razie potrzeby wzywa służby. Brak obrazu nie jest dowodem alarmu fałszywego. Przywrócenie: po fizycznym sprawdzeniu wejścia, ustaleniu przyczyny, naprawie linii i próbie sabotażu; uzbrojenie dopiero po opuszczeniu strefy przez ludzi.

\subsection*{KD, wymuszenie drzwi lub zgubiona karta}
Wyzwalacz: otwarcie bez zezwolenia, zbyt długie otwarcie, uszkodzenie czytnika lub zgłoszenie utraty karty. Odpowiedzialny: administrator KD i dyżurny. System oznacza przejście; patrol sprawdza drzwi i niezależne działanie klamki. Administrator unieważnia kartę i sprawdza synchronizację wszystkich pięciu kontrolerów. Odłączony kontroler wymaga lokalnej aktualizacji albo dozoru zastępczego. Przywrócenie: potwierdzone odebranie uprawnienia, próba odmowy dla utraconej karty, akceptacji właściwej karty/PIN-u i swobodnego wyjścia. Nie stosuje się blokady ewakuacji jako sposobu przeciwdziałania utracie identyfikatora.

\subsection*{Zanik energii i wyczerpanie rezerwy}
Wyzwalacz: komunikat UPS lub centrali o utracie sieci. Odpowiedzialny: dyżurny, serwis elektryczny. Dyżurny sprawdza źródła rezerwowe i zapisuje czas zaniku; nie oczekuje startu nieistniejącego agregatu. Jeśli do 20 min nie ma przewidywanego powrotu energii, zatrzymuje przyjmowanie nowych osób, zabezpiecza materiał i rozpoczyna kontrolowane zamknięcie zapisu/udziału NAS. Najpóźniej po 30 min lub wcześniej przy alarmie niskiej baterii wyłącza komputer i NAS zgodnie z instrukcjami. Utrata PoE kończy obraz i elektryczny dostęp; patrol zapewnia wejścia zastępczo, a klamki pozostają sprawne. Przy zagrożeniu utratą SSP budynek jest opróżniany i pozostaje pod dozorem. Przywrócenie: stabilna sieć, sprawne obwody i RPO, naładowanie do wymaganego poziomu, test zapisu, uprawnień i central; po PWP dodatkowo zgoda na ponowne załączenie i lokalny start UPS.

\subsection*{Utrata LAN, VMS lub NAS}
Wyzwalacz: utrata urządzenia, przerwa zapisu albo brak panelu. Odpowiedzialny: administrator IT i dyżurny. Przy sprawnym PoE kamery przechodzą na SD, KD pracuje lokalnie. Dyżurny korzysta z fizycznych paneli SSP i SSWiN. Awaria SW usuwa również PoE, więc nie zakłada się wtedy dalszego bufora SD. Przywrócenie: kontrola kabli, czasu, poprawnych strumieni i odzyskania luki z SD; próba odtworzenia nagrań z NAS. Nie odwołuje się zastępczego dozoru tylko na podstawie świecącego portu sieciowego.

\subsection*{Cyberincydent i awaria analityki}
Wyzwalacz: nieuprawnione konto, zmiana konfiguracji, podejrzany ruch albo seria fałszywych alarmów AI. Odpowiedzialny: administrator z dyżurnym. Izoluje objęty segment, zachowuje logi i konfigurację, czasowo wyłącza regułę AI1, gdy obciąża ochronę. Nie izoluje fizycznie zasilania SSP ani mechaniki wyjść. Przywrócenie: ustalona przyczyna, odtworzenie zatwierdzonej konfiguracji, zmiana kluczy po podejrzeniu ujawnienia, sprawdzenie ról i ponowne testy scenariuszy. Regułę AI włącza się dopiero po ponownej ocenie FP/FN na oddzielonym zbiorze.

Każdy zapis incydentu obejmuje czas, ID, działania, odpowiedzialnego i warunek zamknięcia. Materiał dowodowy eksportuje uprawniona osoba z kontrolą integralności, np. SHA-256; kopia podlega odrębnym terminom przechowywania.
